% ================================================================
% ==== FILE: content/k_levels/k1.tex
% ================================================================

% ==============================
%  Ontology of Continua — Core
%  K-level Module: K1
%  File: content/k_levels/k1.tex
%  Status: FULLY DEFINED
% ==============================

\subsubsection{\texorpdfstring{$K_1$}{K_1} Übersicht}
\label{sec:k1-overview}

Level $K_1$ ist das erste wirklich \emph{strukturelle} Kontinuum.
Es ergibt sich aus der Wirkung des Operators $\Psi_{0\to1}$ auf $K_0$,
Einführung einer eindimensionalen Achse $A_1$ und damit die Ermöglichung der Existenz
kontinuierlicher Freiheitsgrade.

$K_1$ ist noch kein physischer Raum; vielmehr ist es der abstrakte Prototyp von
ein eindimensionales Kontinuum:
\[
K_1 = (X, \tau, \mu, A_1),
\]
wobei $X$ eine eindimensionale Domäne ist, $\tau$ die Struktur ist
Topologie, und $\mu$ ist ein mit $\tau$ kompatibles Maß.

Die Erstellung von $K_1$ ist der erste Schritt, bei dem:
\begin{itemize}
    \item-Unterschiede können entlang einer Achse ausgedrückt werden,
    \item kohärente Variation wird möglich,
    \item strukturelle Spannungen können sich ausbreiten,
    \item-Einschränkungen erzeugen kontinuierliche Kurven.
\end{itemize}

Jedes höhere Kontinuum erbt $A_1$ als seine Grundachse.

% ================================================================
\subsubsection{Zustandsraum \texorpdfstring{$\Omega(K_1)$}{\Omega(K_1)}}

Der Zustandsraum von $K_1$ ist:
\[
\Omega(K_1) = C^0(X, V),
\]
der Raum stetiger Funktionen aus dem eindimensionalen Bereich $X$
auf einen Wert gesetzt $V$.

Eigenschaften:
\begin{enumerate}
    \item $\Omega(K_1)$ ist unendlichdimensional (Funktionsraum).
    \item Jede Konfiguration ist per Definition kontinuierlich.
    \item Diskontinuitäten liegen außerhalb von $\Omega(K_1)$ und entsprechen
          Verletzungen der strukturellen Kohärenz.
    \item $X$ kann je nach offen, geschlossen, kompakt oder nicht kompakt sein
          Einschränkungen im entsprechenden $M_1$.
\end{enumerate}

Dies ist der minimale Raum, in dem kontinuierliche Variation existiert.

% ================================================================
\subsubsection{Grenze \texorpdfstring{$\partial\Omega(K_1)$}{\partial\Omega(K_1)}}

Die Grenze besteht aus Konfigurationen, die dem Verlust von nahekommen
Kontinuität:
\[
\partial\Omega(K_1) =
\{ f \in C^0(X,V) : \exists\, x_0 \in X \text{ mit }
\lim_{x \to x_0} f(x) \text{ undefiniert} \}.
\]

Interpretation:
\begin{itemize}
    \item Punkte, an denen die Kontinuität fehlschlägt.
    \item Punkte, an denen die strukturelle Spannung $\Theta_1$ überschreitet.
    \item Punkte, an denen die Domäne $X$ degeneriert oder zusammenbricht.
\end{itemize}

Das Überschreiten der Grenze zerstört das Kontinuum $K_1$.

% ================================================================
\subsubsection{Achsen \texorpdfstring{$A(K_1)$}{A(K_1)}}

$K_1$ hat genau eine Achse:
\[
A_1 : X \to \mathbb{R},
\]
stellt den intrinsischen Parameter dar, entlang dessen Unterschiede entstehen
bestellt.

Bedingungen:
\begin{itemize}
    \item $A_1$ muss monoton sein (Axiom A$_{1,\text{mon}}$).
    \item $A_1$ muss zusammenhängend sein (Axiom A$_{1,\text{conn}}$).
    \item $A_1$ muss eine mit $\tau$ kompatible Variation zulassen.
\end{itemize}

Alle zukünftigen Achsen $A_k$ für $k \ge 2$ erweitern diese Grundstruktur.

% ================================================================
\subsubsection{Potentiale \texorpdfstring{$P(K_1)$}{P(K_1)}}

Zu den Potenzialen für $K_1$ gehören:
\[
P(K_1) = \{ P_{\mathrm{grad}}, P_{\mathrm{bound}}, P_{\mathrm{smooth}} \},
\]
beschreibend:
\begin{itemize}
    \item-Verläufe entlang der Achse,
    \item grenzinduzierte Spannungen,
    \item Glättebeschränkungen,
    \item erlaubte die Krümmung von Funktionen in $\Omega(K_1)$.
\end{itemize}

Das sind noch keine Kräfte; es handelt sich um regierende strukturelle Potentiale
Kohärenz der Variation.

% ================================================================
\subsubsection{Schwellenwerte \texorpdfstring{$\Theta(K_1)$}{\Theta(K_1)}}

Schwellenwerte regeln die Zulässigkeit von Zuständen:
\[
\Theta(K_1) = \{ \Theta_{\mathrm{cont}}, \Theta_{\mathrm{grad}},
                 \Theta_{\mathrm{glatt}} \}.
\]

Beispiele:
\begin{itemize}
    \item $\Theta_{\mathrm{cont}}$ – minimale Kontinuität erforderlich.
    \item $\Theta_{\mathrm{grad}}$ — maximal zulässiger Gradient davor
          Spannung wird unhaltbar.
    \item $\Theta_{\mathrm{smooth}}$ – minimale Regelmäßigkeit erforderlich.
\end{itemize}

Wird ein Schwellenwert überschritten, liegt die Konfiguration außerhalb
$\Omega(K_1)$.

% ================================================================
\subsubsection{Flüsse \texorpdfstring{$J(K_1)$}{J(K_1)}}

Flüsse beschreiben erlaubte kontinuierliche Änderungen entlang $A_1$:
\[
J(K_1) = \{ \partial_x f(x),\; \partial_t f(x,t),\; \text{zulässige Verformungen} \}.
\]

Haupteigenschaften:
\begin{itemize}
    \item-Ströme verbreiten strukturelle Spannungen,
    \item-Flüsse müssen die Kontinuität wahren,
    \item Flows können keine neuen Achsen erstellen (erfordert einen Übergang von $K_1\to K_2$).
\end{itemize}

% ================================================================
\subsubsection{Zyklen \texorpdfstring{$C(K_1)$}{C(K_1)}}

Weil $K_1$ kontinuierliche Variation unterstützt, aber noch nicht mehrdimensional
Wechselwirkung, Zyklen sind topologisch trivial:
\[
C_{\mathrm{triv}} = \text{konstante Schleifen im Funktionenraum}.
\]

Es gibt keine nichttrivialen Strukturzyklen, da die zyklische Struktur dies erfordert
mindestens zwei unabhängige Achsen ($K_2$).

Also:
\[
\tau(K_1) \text{ kann nicht entstehen.}
\]

% ================================================================
\subsubsection{Zeit \texorpdfstring{$\tau(K_1)$}{\tau(K_1)}}

Bei $K_1$ existiert keine Zeit.

Nach dem Time-Origin-Theorem (Physics Run):

\[
\text{Zeit entsteht erst bei } K_2 \text{ aus nicht abbaubaren Zyklen}.
\]

Also:
\[
\tau(K_1) \text{ ist undefiniert und nicht darstellbar.}
\]

% ================================================================
\subsubsection{Kontinuität \texorpdfstring{$k(K_1)$}{k(K_1)}}

Das Maß der Kontinuität:
\[
k(K_1) =
H(\Omega(K_1)) \cdot
\frac{|A(K_1)|}{|A(M_1)|} \cdot
(1 - \frac{T_1}{\Theta_{\mathrm{cont}}})_+,
\]
wo:
\begin{itemize}
    \item $|A(K_1)| = 1$,
    \item $A(M_1)$ erlaubt möglicherweise mehr Achsen, ist aber noch nicht realisiert,
    \item $T_1$ ist strukturelle Spannung,
    \item $(\cdot)_+$ ist der positive Teil.
\end{itemize}

$k(K_1)$ misst, ob $K_1$ strukturell tragfähig bleibt.

% ================================================================
\subsubsection{Strukturelle Spannung \texorpdfstring{$T(K_1)$}{T(K_1)}}

Strukturelle Spannung auf $K_1$ entsteht durch:
\begin{itemize}
    \item Gradienten von Funktionen,
    \item Grenzeffekte,
    \item Inkompatibilität von Einschränkungen entlang $A_1$.
\end{itemize}

Wenn:
\[
T(K_1) > \Theta_{\mathrm{cont}},
\]
dann versagt die Kontinuität und das Kontinuum bricht zusammen.

% ================================================================
\subsubsection{Energie \texorpdfstring{$E(K_1)$}{E(K_1)}}

Energie wird strukturell definiert als:
\[
E(K_1) = \int_X \mathcal{E}(f(x), \partial_x f(x))\, d\mu,
\]
wobei $\mathcal{E}$ eine strukturelle Energiedichte ist, abhängig von:
\begin{itemize}
    \item Wertvariation,
    \item Glättebeschränkungen,
    \item Regelmäßigkeitsstrafen.
\end{itemize}

$E(K_1)$ hat noch keine physikalische Bedeutung; es handelt sich um eine funktionale Messung
Kohärenz.

% ================================================================
\subsubsection{Operatoren auf \texorpdfstring{$K_1$}{K_1} (\texorpdfstring{$\Psi$}{\Psi}, \texorpdfstring{$\Phi$}{\Phi}, \texorpdfstring{$\Lambda$}{\Lambda}, \texorpdfstring{$U$}{U}, \texorpdfstring{$\Chi$}{\Chi})}

\begin{itemize}
    \item $\Psi_{1\to2}$ – Geburt der zweiten Achse $A_2$, Aktivierung
          $K_2$.
    \item $\Phi$ – Entwicklung des zulässigen Metaraums $M_1$.
    \item $\Lambda$ – Einschränkungsoperator, der Kontinuität erzwingt,
          Monotonie und Verbundenheit.
    \item $U$ – Vereinheitlichung der lokalen Struktur in globale Kohärenz.
    \item $\Chi$ – Verzweigungsoperator, der inkompatibles $K_1$ ermöglicht
          Mannigfaltigkeiten.
\end{itemize}

Diese Operatoren definieren, wie sich $K_1$ entwickelt und wie neue Dimensionen entstehen.

% ================================================================
\subsubsection{Prozesse auf \texorpdfstring{$K_1$}{K_1}}

Schlüsselprozesse:
\begin{itemize}
    \item Ausbreitung von Gradienten,
    \item Glättung von Unstetigkeiten,
    \item Stabilisierung kontinuierlicher Konfigurationen,
    \item struktureller Zusammenbruch unter übermäßiger Spannung,
    \item Dimensionsübergang zu $K_2$ unter Druck von
          inkompatible Unterschiede (Axiom~15).
\end{itemize}

% ================================================================
\subsubsection{Vorhersagen für \texorpdfstring{$K_1$}{K_1}}

Die Theorie sagt voraus:
\begin{enumerate}
    Die \item-Kontinuität muss global gewahrt bleiben.
    Das Dimensionswachstum von \item erfordert inkompatible Variationen, die nicht möglich sind
          auf einer einzigen Achse dargestellt werden;
    \item bei $K_1$ treten keine zeitlichen Phänomene auf;
    \item-Kollaps unter übermäßigen Farbverläufen ist universell;
    \item alle $K_1$ Kontinuen müssen in ihren $M_1$ Metaraum eingebettet werden.
\end{enumerate}

% ================================================================
\subsubsection{Experimente für \texorpdfstring{$K_1$}{K_1}}

Indirekte Tests:
\begin{itemize}
    \item Universalität des 1D-Gradientenkollapsverhaltens in physikalischen,
          chemische und biologische Filamente,
    \item Beobachtung des Schwellenwertverhaltens bei steilen Gefällen,
    \item Erkennung von erzwungenem Dimensionswachstum in Systemen, in denen a
          Eine einzelne Achse wird unzureichend.
\end{itemize}

Beispiele:
\begin{itemize}
    \item Polymerstrang kollabiert unter Spannung,
    \item memristive Filamentzerstörung,
    \item 1D-Reaktions-Diffusionskanal-Instabilität.
\end{itemize}

% ================================================================
\subsubsection{Zusammenbruch und Tod von \texorpdfstring{$K_1$}{K_1}}

$K_1$ bricht zusammen, wenn:
\[
\Omega(K_1) = \varnothing.
\]

Dies geschieht, wenn:
\begin{itemize}
    Die \item-Kontinuität ist unterbrochen,
    \item-Gradienten überschreiten Schwellenwerte,
    \item-Einschränkungen in $\mathcal{L}$ (übertragen von $K_0$) werden
          inkompatibel.
\end{itemize}

Nach dem Zusammenbruch ist der Übergang zu $K_2$ unmöglich.

% ================================================================
\subsubsection{Falschbarkeit von \texorpdfstring{$K_1$}{K_1}}

Vorhersagen, die empirisch verfälscht werden können:
\begin{enumerate}
    \item universelle Präsenz von Kontinuitätsschwellen,
    \item Fehlen nichttrivialer Zyklen,
    \item Fehlen entstehender Zeit,
    \item Anforderung an monotone verbundene Achsen,
    \item-Dimensionsübergang nur bei inkompatiblen Unterschieden.
\end{enumerate}

Ein Verstoß gegen eine dieser Vorhersagen stellt die $K_1$-Beschreibung in Frage.

% ================================================================
\subsubsection{Verzweigung / ontologische Position von \texorpdfstring{$K_1$}{K_1}}

$K_1$ liegt direkt über $K_0$ im ontologischen Baum:
\[
K_0 \to K_1 \to K_2.
\]

Die Verzweigung bei $K_1$ entspricht inkompatiblen zulässigen Mannigfaltigkeiten
Sie teilen die gleichen Metagesetze, unterscheiden sich jedoch in den strukturellen Einschränkungen.

Solche Verzweigungen bleiben zulässig, solange sie innerhalb von $\Omega(K_1)$ bleiben.

% ================================================================
\subsubsection{Beziehung zu M-Räumen}

$K_1$ ist in seinen Metaraum eingebettet:
\[
K_1 \subseteq M_1 \subseteq \Omega(K_0).
\]

$M_1$ bestimmt:
\begin{itemize}
    \item die Topologie $\tau$ von $X$,
    \item die Regularitätsbeschränkungen,
    \item zulässige Bereiche für Potentiale,
    \item Zulässigkeitsschwellen.
\end{itemize}

Der Übergang zu $K_2$ erfolgt, wenn:
\[
A_1 \text{ kann nicht alle auftretenden Unterschiede darstellen.}
\]
